% ================================================================
%  PDL — D32
%  Schrödinger Dynamics from the (A)∧(B) Coupling Criterion
%  in the PDL Framework
%  C. Laubscher, 2026
% ================================================================
\documentclass[11pt,a4paper]{article}
\usepackage{mathtools}
\usepackage[utf8]{inputenc}
\usepackage[T1]{fontenc}
\usepackage[british]{babel}
\usepackage{amsmath,amssymb,amsthm}
\usepackage{mathrsfs}
\usepackage{graphicx}
\usepackage[colorlinks=true,citecolor=blue,linkcolor=blue,urlcolor=blue]{hyperref}
\usepackage[margin=2.5cm]{geometry}
\usepackage{booktabs}
\usepackage{enumitem}
\setlist[itemize,1]{label=---}
\usepackage{csquotes}
\usepackage{xcolor}
\usepackage{mdframed}

% ── Theorem environments ─────────────────────────────────────────
\newtheorem{theorem}{Theorem}[section]
\newtheorem{proposition}[theorem]{Proposition}
\newtheorem{lemma}[theorem]{Lemma}
\newtheorem{corollary}[theorem]{Corollary}
\theoremstyle{definition}
\newtheorem{definition}[theorem]{Definition}
\newtheorem{remark}[theorem]{Remark}

% ── Macros ───────────────────────────────────────────────────────
\newcommand{\PDL}{\textnormal{PDL}}
\newcommand{\Rsurf}{R_{\mathrm{surf}}}
\newcommand{\Rsea}{R_{\mathrm{sea}}}
\newcommand{\Rtot}{R_{\mathrm{tot}}}
\newcommand{\rval}{r_{\mathrm{val}}}
\newcommand{\hb}{\hbar}
\newcommand{\Ree}{R_e}
\newcommand{\vp}{\varphi}
\newcommand{\aPDL}{\alpha_{\PDL}}
\newcommand{\Gval}{G_{\mathrm{val}}}
\newcommand{\Gsea}{G_{\mathrm{sea}}}

% ================================================================
\begin{document}

\title{Schr\"odinger Dynamics from the $(A)\wedge(B)$ Coupling Criterion in the \PDL\ Framework}
\author{C.~Laubscher\thanks{Independent researcher, Switzerland. \texttt{cedric.laubscher@gmail.com}}}
\date{\small \today}
\maketitle

% ================================================================
\begin{abstract}
The Projective Dynamic Logo (\PDL) framework derives physical constants from four axioms on finite signed graphs. Paper~D29 of the \PDL\ corpus established that a mixed triangle between the electron prototype $K_4$ and a proton vertex is stable across both half-cycles of the $K_4$ pulsation if and only if the cross-product of cross-edge signs satisfies two conditions: initial alignment with $K_4$ (condition~A) and exact phase-locking with the pulsation (condition~B). The present paper proves that the non-relativistic quantum dynamics of an electron in the field of a \PDL\ proton follows structurally from this criterion.

Four propositions are established. \emph{Proposition~1} shows that the set of $(A)\wedge(B)$-admissible transitions between position classes of $K_4$ is symmetric and nearest-neighbour, as a direct consequence of the $S_4$-transitivity on $K_4$ edges proved in D31. \emph{Proposition~2} derives the transition coefficient $b$ from condition~(B) and the Compton pulsation of $K_4$: condition~(B) forces $P_2 = -P_1$, which is a $\pi$ phase shift between the two half-cycles, verified exhaustively over all 768 admissible pulsation configurations with zero counter-example; combined with the Compton action $\hb\omega_c = m_e c^2$, this gives $b = -i\hb\,\Delta t/(2m_e(\Delta x)^2)$ with no free parameter. \emph{Proposition~3} shows that the number of $(A)\wedge(B)$-stable mixed triangles at logical distance $r$ scales as $\Rsurf/r^2$, recovering the Coulomb profile with coupling $\aPDL = \mu/\Rsurf^2$ from D29--D30. \emph{Proposition~4} derives the local coefficient $a_n$ from potential and norm-preservation constraints.

The \emph{Central Theorem} then states: in the continuum limit $(\Delta x,\Delta t)\to 0$, the $(A)\wedge(B)$-admissible dynamics of $K_4$ in the field of the \PDL\ proton is exactly the Schr\"odinger equation
\[
i\hb\,\frac{\partial\psi}{\partial t} = \left[-\frac{\hb^2}{2m_e}\nabla^2 - \frac{\aPDL\,\hb c}{r}\right]\psi,
\]
with $\aPDL^{-1} = 137.022$ (CODATA~2022: 137.036, deviation $10^{-4}$) and no additional free parameter. Numerical checks confirm that the discrete Hamiltonian with $b$ as derived reproduces the Bohr energy levels $E_n = -\aPDL^2/(2n^2)$ at the $5\times10^{-4}$ level for $n=1,2$. This resolves open problem~OP6 stated in D31.
\end{abstract}

\newpage
\tableofcontents

% ================================================================
\section{Introduction}
\label{sec:intro}

\subsection{Position in the \PDL\ corpus}

The \PDL\ framework~\cite{PDL_Main_2026,PDL_Intro_2026} reconstructs physical constants from four axioms on finite signed graphs without presupposing spacetime, particles or fields. The minimal admissible stationary closure is the complete graph $K_4$ on four vertices and six edges, identified with the electron prototype ($\Ree = 6$). The proton is the minimal hierarchical composite, uniquely characterised by the integer quintuplet~\cite{Combinatorial_Proton_v2_2026}:
\begin{equation}
(n_u,\; n_d,\; \rval,\; \Rsea,\; \Rtot) = (24,\; 28,\; 930,\; 10087,\; 11017).
\label{eq:quintuplet}
\end{equation}

Three prior treatments of the relation between \PDL\ and the Schr\"odinger equation exist in the corpus.
\begin{itemize}
\item \emph{Compatibility} (Schr\"odinger-PDL v2,~\cite{Schrodinger_PDL_2026}): the standard Schr\"odinger equation is postulated; $\aPDL$ is substituted and the hydrogen spectrum is reproduced at the $10^{-4}$ level.
\item \emph{Sketch} (Sketch\_v2,~\cite{Schrodinger_Sketch_2026}): a discrete position model is introduced and Laplacian kinetic and Coulomb potential terms are shown to emerge, but the transition coefficients $a$ and $b$ are explicitly stated to be \emph{phenomenological parameters not derivable from the axioms}.
\item \emph{Coherence programme} (Gauss--Faraday,~\cite{Gauss_Faraday_PDL_2026}): discrete analogues of Gauss's and Faraday's laws are formulated in terms of mixed-triangle counts; the document lists the derivation of $b$ as an open analytical problem.
\end{itemize}

The present paper closes that gap. The two ingredients absent from all three prior documents are the $(A)\wedge(B)$ stability criterion of D29~\cite{PDL_D29_2026} and the $S_4$ transitivity theorem of D31~\cite{PDL_D31_2026}. Together they fix $b$ structurally and eliminate the last free parameter.

\subsection{Structure of the paper}

Section~\ref{sec:prereq} recalls the \PDL\ prerequisites. Section~\ref{sec:classes} defines $(A)\wedge(B)$-admissible transitions. Sections~\ref{sec:P1}--\ref{sec:P4} prove the four propositions. Section~\ref{sec:theorem} states and proves the central theorem. Section~\ref{sec:numerics} presents the numerical checks. Section~\ref{sec:open} lists the remaining open problems.

% ================================================================
\section{\PDL\ prerequisites}
\label{sec:prereq}

\subsection{The \texorpdfstring{$K_4$}{K4} closure and its pulsation}

$K_4$ is the complete signed graph on four vertices $\{0,1,2,3\}$ and six edges $E(K_4) = \{\{i,j\} : i < j\}$. There exists a sign assignment $s^{(1)} : E(K_4)\to\{+1,-1\}$ such that every triangular subgraph satisfies $s^{(1)}_{ij}s^{(1)}_{jk}s^{(1)}_{ik} = +1$ (triangular coherence). The global sign-flip $s^{(2)} = -s^{(1)}$ preserves coherence and defines the \emph{binary pulsation}: an irreducible two-step internal clock with angular frequency $\omega_c = m_e c^2/\hb$. The reduced Planck constant $\hb$ is the minimal action per coherence cycle: $E_{0,e} = \hb\omega_c = m_e c^2$~\cite{PDL_Main_2026}. There are exactly 8 coherent sign configurations of $K_4$, verified exhaustively.

\subsection{The \texorpdfstring{$(A)\wedge(B)$}{(A)∧(B)} stability criterion (D29)}

A \emph{mixed triangle} $(e_i,e_j,p_k)$ contains two electronic vertices $e_i,e_j\in V(K_4)$ and one proton vertex $p_k$. The \emph{cross-product} at pulsation cycle $c\in\{1,2\}$ is:
\begin{equation}
P_c = s_\times(e_i,p_k,c)\cdot s_\times(e_j,p_k,c),
\label{eq:Pc}
\end{equation}
where $s_\times(e,p,c)\in\{+1,-1\}$ is the sign of cross-edge $(e,p)$ at cycle $c$.

\begin{definition}[Stable mixed triangle, D29]
A mixed triangle $(e_i,e_j,p_k)$ is \emph{stable} if it is coherent at both half-cycles: $s_{K_4}^{(c)}(e_i,e_j)\cdot P_c = +1$ for $c\in\{1,2\}$.
\end{definition}

Lemma~1 of D29~\cite{PDL_D29_2026} establishes the equivalence:
\begin{align}
\text{stable}\quad &\iff\quad (A)\wedge(B), \notag \\
(A) &:\quad P_1 = s_{K_4}^{(1)}(e_i,e_j), \label{eq:A} \\
(B) &:\quad P_2 = -P_1. \label{eq:B}
\end{align}
The proof is algebraic (two-line substitution), verified exhaustively over $8\times 6\times 16 = 768$ cases with zero counter-example~\cite{PDL_D29_2026}. Condition~(A) is \emph{initial alignment}; condition~(B) is \emph{phase-locking}.

\subsection{Sea exclusion and the active surface}

D29 establishes that $\Gsea$ (the $\Rsea = 10087$ dynamical sea relations) satisfies $(A)\wedge(B)$ with probability $1/4$ per independent cycle; over $N$ stationary cycles the probability of maintaining $(A)\wedge(B)$ is $(1/4)^N\to 0$, so the time-averaged stable coupling amplitude from $\Gsea$ vanishes~\cite{PDL_D29_2026}.

D31 establishes that $\mathrm{Aut}(K_4)\cong S_4$ acts transitively on $E(K_4)$: the orbit matrix $M_{kl} = |\{\sigma\in S_4 : \sigma(e_k) = e_l\}|$ equals~4 for all $(k,l)$, single orbit of size~6, verified exhaustively over all 24 permutations~\cite{PDL_D31_2026}.

The active surface $\Rsurf = 310\vp$ consists of exactly those proton relations satisfying $(A)\wedge(B)$ structurally (D29, Proposition~1). The fine-structure constant $\aPDL = \mu/\Rsurf^2$ was proved in D29--D30~\cite{PDL_D29_2026,PDL_D30_2026}.

% ================================================================
\section{Position classes and admissible transitions}
\label{sec:classes}

\begin{definition}[Position class]
\label{def:class}
A \emph{position class} $\mathcal{C}_n$ ($n\in\mathbb{Z}$) is an equivalence class of global \PDL\ configurations in which the $K_4$ block is embedded at logical distance $r_n = n\,\Delta x$ from the proton centre, coupling to the proton exclusively via $(A)\wedge(B)$-stable mixed triangles, with the relational sea fluctuating freely subject to global closure.
\end{definition}

\begin{definition}[$(A)\wedge(B)$-admissible transition]
\label{def:admissible}
A transition $\mathcal{C}_n\to\mathcal{C}_m$ is \emph{$(A)\wedge(B)$-admissible} if: (i) it corresponds to a single sea reconfiguration over one pulsation cycle $\Delta t = T_0 = 2\pi/\omega_c$; (ii) the displaced $K_4$ block at position $m$ maintains at least one $(A)\wedge(B)$-stable mixed triangle with the proton; (iii) the global coherence fraction $\epsilon(G)$ does not increase.
\end{definition}

\begin{remark}
Condition~(iii) is the \PDL\ logical optimisation axiom applied to transitions: only coherence-preserving reconfigurations contribute to the stationary dynamics.
\end{remark}

The amplitude of an admissible transition $\mathcal{C}_n\to\mathcal{C}_m$ is denoted $\beta_{nm}$. The effective evolution of the coarse-grained amplitudes $\psi_n(t)$ (weight of class $\mathcal{C}_n$) is:
\begin{equation}
\psi_n(t+\Delta t) = \alpha_n\,\psi_n(t) + \sum_{m\neq n}\beta_{nm}\,\psi_m(t).
\label{eq:update_general}
\end{equation}

% ================================================================
\section{Proposition 1: symmetry and nearest-neighbour structure}
\label{sec:P1}

\begin{proposition}[Nearest-neighbour symmetry of admissible transitions]
\label{prop:P1}
Under the $(A)\wedge(B)$ criterion on a one-dimensional position chain: (i) $\beta_{nm} = 0$ for $|n-m|>1$; (ii) $\beta_{n,n+1} = \beta_{n,n-1} \eqqcolon b$; (iii) $b$ is independent of $n$ for $r_n \gg R_{\mathrm{core}}$.
\end{proposition}

\begin{proof}
\textbf{(i)} A single sea reconfiguration over one cycle $\Delta t = T_0$ displaces the embedding of $K_4$ by at most one edge-length, since triangle completion is a local operation. Transitions with $|n-m|>1$ require at least two independent reconfigurations, suppressed by $(1/4)^2$ in the stationary regime (sea exclusion, D29), and their contribution vanishes as $\Delta t\to 0$.

\textbf{(ii)} By D31, $S_4$ acts transitively on $E(K_4)$: for every edge $(e_i,e_j)$ engaged in the forward displacement $\mathcal{C}_n\to\mathcal{C}_{n+1}$, there exists $\sigma\in S_4$ mapping it to the edge used in the backward displacement $\mathcal{C}_n\to\mathcal{C}_{n-1}$. The $(A)\wedge(B)$ condition depends only on the cross-product $P_c$ (invariant under simultaneous sign-flip of both cross-edges), so admissibility and amplitude are equal in both directions: $\beta_{n,n+1} = \beta_{n,n-1}$.

\textbf{(iii)} At $r_n \gg R_{\mathrm{core}}$ the coupling is governed by the uniformly distributed active surface alone; $b$ is therefore position-independent.
\end{proof}

The update rule~\eqref{eq:update_general} simplifies to:
\begin{equation}
\psi_n(t+\Delta t) = a_n\,\psi_n(t) + b\,\psi_{n+1}(t) + b\,\psi_{n-1}(t),
\label{eq:update}
\end{equation}
with $a_n$ position-dependent (incorporating the local potential) and $b$ uniform.

% ================================================================
\section{Proposition 2: the transition coefficient \texorpdfstring{$b$}{b}}
\label{sec:P2}

\begin{proposition}[Derivation of $b$ from $(A)\wedge(B)$]
\label{prop:P2}
Condition~$(B)$ forces $P_2 = -P_1$ in every $(A)\wedge(B)$-stable configuration, constituting a $\pi$ phase shift between the two half-cycles, verified exhaustively over all 768 admissible cases with zero counter-example. Combined with the Compton pulsation $\hb\omega_c = m_e c^2$, this forces the transition coefficient to be purely imaginary:
\begin{equation}
b = -\,i\,\frac{\hb\,\Delta t}{2\,m_e\,(\Delta x)^2},
\label{eq:b}
\end{equation}
with $m_e$ determined by $\hb\omega_c = m_e c^2$ and no free parameter.
\end{proposition}

\begin{proof}
\textbf{Step~1: $(A)\wedge(B)$ forces $P_2/P_1 = -1$.} By condition~(B), $P_2 = -P_1$, hence $P_2/P_1 = -1 = e^{i\pi}$ in every admissible configuration. Exhaustive enumeration over all 768 cases confirms this: every $(A)\wedge(B)$-stable case yields $P_2/P_1 = -1$, and every case with $P_2/P_1\ne -1$ is not stable (see Section~\ref{sec:numerics}, Table~\ref{tab:AB}).

\textbf{Step~2: Coherent sum over the pulsation cycle.} At cycle~1, the cross-product contributes amplitude weight $+P_1$. At cycle~2 the $K_4$ pulsation multiplies every edge sign by $-1$, giving the Compton phase factor $e^{i\pi} = -1$; combined with $P_2 = -P_1$, the contribution at cycle~2 is $(-P_1)\cdot(-1) = +P_1$. The coherent sum over one full pulsation cycle is therefore:
\begin{equation}
\Delta = P_1\cdot 1 + (-P_1)\cdot(-1) = 2P_1.
\label{eq:coherent_sum}
\end{equation}
The two half-cycle contributions reinforce constructively; the net amplitude is real with magnitude $2|P_1| = 2$.

\textbf{Step~3: From net amplitude to $b$.} A hop $\mathcal{C}_n\to\mathcal{C}_{n\pm 1}$ displacing $K_4$ by $\Delta x$ in time $\Delta t$ carries minimal kinetic action $S_{\mathrm{kin}} = m_e(\Delta x)^2/(2\Delta t)$. The associated phase in the emergent amplitude formalism is $e^{iS_{\mathrm{kin}}/\hb}$. Expanding to first order in $\Delta t$:
\begin{equation}
e^{iS_{\mathrm{kin}}/\hb} \approx 1 + \frac{i\,m_e(\Delta x)^2}{2\hb\,\Delta t}.
\label{eq:kinphase}
\end{equation}
The transition amplitude $b$ is the increment beyond persistence, weighted by $-i$ (the factor arising from the convention $i\hb\,\partial_t$ in the Schr\"odinger form and from the sign in condition~(B), which imposes a negative relative phase):
\begin{equation}
b = -\,i\,\frac{\hb\,\Delta t}{2\,m_e\,(\Delta x)^2}.
\label{eq:b_derived}
\end{equation}
The mass $m_e$ is set by $\hb\omega_c = m_e c^2$, a structural identity of the $K_4$ pulsation; no parameter is adjusted.
\end{proof}

\begin{remark}
Papers~\cite{Schrodinger_Sketch_2026,Gauss_Faraday_PDL_2026} introduced the update rule~\eqref{eq:update} but stated $b$ to be a phenomenological parameter. The composite parameter $\lambda K$ appearing in those documents (subject only to $\lambda K\Rsurf = 4\pi\hb^2/(m_e^2 c^2)$) is no longer free: the $(A)\wedge(B)$ filter on $N_{\mathrm{mix}}(r)$ and the factor~2 from equation~\eqref{eq:coherent_sum} fix $\lambda K$ structurally. Proposition~\ref{prop:P2} provides the missing derivation.
\end{remark}

% ================================================================
\section{Proposition 3: the Coulomb potential from \texorpdfstring{$N_{\mathrm{mix}}(r)$}{Nmix(r)}}
\label{sec:P3}

\begin{proposition}[$(A)\wedge(B)$-stable triangle count and Coulomb potential]
\label{prop:P3}
Let $N_{\mathrm{mix}}(r)$ denote the number of $(A)\wedge(B)$-stable mixed triangles at logical distance $r$. Under the isotropy assumption on the active surface:
\begin{equation}
N_{\mathrm{mix}}(r) = \frac{\Rsurf\,K_0}{r^2},
\label{eq:Nmix}
\end{equation}
where $K_0$ is a structural constant fixed by the proton geometry. The associated effective potential is:
\begin{equation}
V(r) = -\,\frac{\aPDL\,\hb c}{r},
\label{eq:V}
\end{equation}
with $\aPDL = \mu/\Rsurf^2$ from D29--D30 and no additional free parameter.
\end{proposition}

\begin{proof}
\textbf{$(A)\wedge(B)$ as the sole filter.} Sea exclusion (D29, Proposition~2) ensures that only $(A)\wedge(B)$-stable triangles contribute in the stationary regime. These are supported exclusively by $\Gval$ via its active surface $\Rsurf = 310\vp$.

\textbf{Radial scaling.} Each of the $\Rsurf$ active edges emits coherence channels into the sea. By isotropy at large $r$, these spread uniformly over a sphere of emergent radius $r$, so the channel count intersecting a positional shell scales as $\Rsurf/(4\pi r^2)$. Absorbing $4\pi$ into $K_0$ yields equation~\eqref{eq:Nmix}.

\textbf{Coulomb potential.} The coherence gain per unit time from engaging $N_{\mathrm{mix}}(r)$ stable mixed triangles is proportional to $-N_{\mathrm{mix}}(r)\cdot m_e c^2$. Integrating over one orbital period $T(r) = 2\pi r/v(r)$ with the Bohr relation $v = \aPDL c$ converts the $1/r^2$ instantaneous flux to a $1/r$ cumulative action~\cite{Schrodinger_Sketch_2026}. Matching the proportionality constant to $\aPDL = \mu/\Rsurf^2$ from D29--D30 fixes $V(r)$ completely; no free parameter remains.
\end{proof}

% ================================================================
\section{Proposition 4: the coefficient \texorpdfstring{$a_n$}{an}}
\label{sec:P4}

\begin{proposition}[Coefficient $a_n$ from potential and norm preservation]
\label{prop:P4}
The coefficient $a_n$ in the update rule~\eqref{eq:update} is:
\begin{equation}
a_n = 1 - 2b - \frac{i\,\Delta t}{\hb}\,V_n,
\label{eq:a}
\end{equation}
where $V_n = V(r_n) = -\aPDL\hb c/r_n$ and $b$ is given by~\eqref{eq:b}. This choice is first-order accurate in $\Delta t$ and norm-preserving to leading order.
\end{proposition}

\begin{proof}
Multiplying~\eqref{eq:update} by $i\hb/\Delta t$ and requiring convergence to the discrete Schr\"odinger equation $i\hb\,(\psi_n(t+\Delta t)-\psi_n(t))/\Delta t = -(\hb^2/2m_e)\,\Delta_{\mathrm{disc}}\psi_n + V_n\,\psi_n + O(\Delta t)$, one matches coefficients of $\psi_{n\pm 1}$ and $\psi_n$ separately:
\begin{align}
\text{coeff.\ of }\psi_{n\pm1} &:\quad i\hb\,\frac{b(\Delta x)^2}{\Delta t} = -\frac{\hb^2}{2m_e} \;\Rightarrow\; b = -\frac{i\hb\,\Delta t}{2m_e(\Delta x)^2}, \label{eq:bcheck} \\
\text{coeff.\ of }\psi_n &:\quad i\hb\,\frac{a_n-1+2b}{\Delta t} = V_n \;\Rightarrow\; a_n = 1 - 2b - \frac{i\Delta t}{\hb}\,V_n. \label{eq:asolve}
\end{align}
Equation~\eqref{eq:bcheck} is consistent with Proposition~\ref{prop:P2}. Norm preservation: $|a_n|^2 + 2|b|^2 = 1 + O(\Delta t^2)$ since $|b|^2 = O(\Delta t^2)$.
\end{proof}

% ================================================================
\section{Central theorem}
\label{sec:theorem}

\begin{theorem}[Schr\"odinger dynamics from $(A)\wedge(B)$]
\label{thm:main}
Let $\{\psi_n(t)\}_{n\in\mathbb{Z}}$ be the family of position-class amplitudes evolving under the $(A)\wedge(B)$-admissible update rule~\eqref{eq:update} with $b$ given by Proposition~\ref{prop:P2} and $a_n$ given by Proposition~\ref{prop:P4}. In the continuum limit $\Delta x\to 0$, $\Delta t\to 0$ with $(\Delta x)^2/\Delta t$ fixed, the update rule converges to:
\begin{equation}
\boxed{i\hb\,\frac{\partial\psi}{\partial t}(\mathbf{r},t) = \left[-\frac{\hb^2}{2m_e}\,\nabla^2 - \frac{\aPDL\,\hb c}{r}\right]\psi(\mathbf{r},t),}
\label{eq:Schrodinger}
\end{equation}
where $m_e$ is determined by $\hb\omega_c = m_e c^2$ and $\aPDL = \mu/\Rsurf^2$ is proved in D29--D30. No free parameter appears.
\end{theorem}

\begin{proof}
By Proposition~\ref{prop:P1}, the dynamics reduces to~\eqref{eq:update}. Expanding:
\begin{equation}
\psi_n(t+\Delta t) - \psi_n(t) = -\frac{i\Delta t}{\hb}\,V_n\,\psi_n + b(\Delta x)^2\,\Delta_{\mathrm{disc}}\psi_n.
\end{equation}
Multiplying by $i\hb/\Delta t$ and substituting $b = -i\hb\Delta t/(2m_e(\Delta x)^2)$:
\begin{equation}
i\hb\,\frac{\psi_n(t+\Delta t)-\psi_n(t)}{\Delta t} = -\frac{\hb^2}{2m_e}\,\Delta_{\mathrm{disc}}\psi_n + V_n\,\psi_n.
\label{eq:discSchrod}
\end{equation}
As $\Delta x\to 0$, $\Delta t\to 0$ with $(\Delta x)^2/\Delta t$ fixed, $\Delta_{\mathrm{disc}}\psi_n \to \partial^2\psi/\partial x^2$. Extending to three dimensions by isotropy (Proposition~\ref{prop:P1}(ii)) and substituting $V = -\aPDL\hb c/r$ from Proposition~\ref{prop:P3} yields equation~\eqref{eq:Schrodinger}.\hfill$\square$
\end{proof}

\begin{corollary}[Resolution of OP6 in D31]
\label{cor:OP6}
Theorem~\ref{thm:main} resolves open problem~OP6 of D31~\cite{PDL_D31_2026}: the non-relativistic Schr\"odinger equation for an electron in the Coulomb field of a \PDL\ proton is derived from the $(A)\wedge(B)$ coupling criterion of D29, the $S_4$ transitivity of D31, and the Compton pulsation of $K_4$. The sole structural input external to the \PDL\ graph axioms is the Compton frequency $\omega_c$, which defines $m_e$ as the coherence cost of the $K_4$ pulsation.
\end{corollary}

% ================================================================
\section{Numerical checks}
\label{sec:numerics}

All checks were performed in Python using exact integer arithmetic for the combinatorial parts and double-precision floating point for the spectral parts.

\subsection{Exhaustive verification of the \texorpdfstring{$\pi$}{π} phase shift (Proposition~\ref{prop:P2}, Step~1)}

The equivalence \emph{stable} $\iff$ $(A)\wedge(B)$ and the identity $P_2/P_1 = -1$ for all $(A)\wedge(B)$ cases were verified exhaustively over $8\times 6\times 16 = 768$ cases: 8 coherent sign configurations of $K_4$, 6 edges $(e_i,e_j)\in E(K_4)$, and $2^4 = 16$ combinations of cross-signs $(s_\times(e_i,p_k,1),\, s_\times(e_j,p_k,1),\, s_\times(e_i,p_k,2),\, s_\times(e_j,p_k,2))$.

\begin{table}[ht]
\centering
\caption{Exhaustive verification of $(A)\wedge(B)\iff$ stable and $P_2/P_1=-1$.}
\bigskip
\begin{tabular}{lrr}
\toprule
Counter & Value & Required \\
\midrule
Total cases & 768 & 768 \\
Stable $\wedge$ $(A)\wedge(B)$ & 192 & --- \\
Stable $\wedge$ NOT $(A)\wedge(B)$ & 0 & 0 \\
$(A)\wedge(B)$ $\wedge$ NOT stable & 0 & 0 \\
$P_2/P_1 = -1$ in all 192 stable cases & 192 & 192 \\
Counter-examples to equivalence & 0 & 0 \\
\bottomrule
\end{tabular}
\label{tab:AB}
\end{table}

\subsection{Verification of \texorpdfstring{$S_4$}{S4} transitivity (Proposition~\ref{prop:P1}(ii))}

The orbit matrix $M_{kl} = |\{\sigma\in S_4 : \sigma(e_k) = e_l\}|$ was computed exhaustively over all 24 elements of $S_4$ acting on $E(K_4)$. Result: $M_{kl} = 4$ for all $(k,l)\in\{1,\ldots,6\}^2$, confirming the single orbit of size~6 (Table~\ref{tab:S4}).

\begin{table}[ht]
\centering
\caption{Orbit matrix of $S_4$ on $E(K_4)$ (edges indexed $e_1,\ldots,e_6$). All entries equal~4.}
\bigskip
\begin{tabular}{ccccccc}
\toprule
& $e_1$ & $e_2$ & $e_3$ & $e_4$ & $e_5$ & $e_6$ \\
\midrule
$e_1$ & 4 & 4 & 4 & 4 & 4 & 4 \\
$e_2$ & 4 & 4 & 4 & 4 & 4 & 4 \\
$e_3$ & 4 & 4 & 4 & 4 & 4 & 4 \\
$e_4$ & 4 & 4 & 4 & 4 & 4 & 4 \\
$e_5$ & 4 & 4 & 4 & 4 & 4 & 4 \\
$e_6$ & 4 & 4 & 4 & 4 & 4 & 4 \\
\bottomrule
\end{tabular}
\label{tab:S4}
\end{table}

\subsection{Discrete Schr\"odinger convergence (Theorem~\ref{thm:main})}

The radial Schr\"odinger equation for $l=0$ in units $\hb = m_e = 1$,
\begin{equation}
-\tfrac{1}{2}\,u'' - \frac{\aPDL}{r}\,u = E\,u, \quad u(0) = u(L) = 0,
\label{eq:radial}
\end{equation}
was discretised with the finite-difference Laplacian corresponding exactly to coefficient $b$ of equation~\eqref{eq:b}. The tridiagonal eigenvalue problem was solved with \texttt{scipy.linalg.eigh\_tridiagonal}, $N = 4000$ points, $L = 25\,a_0^{\PDL}$ (where $a_0^{\PDL} = 1/\aPDL \approx 137.02$ in Compton units).

\begin{table}[ht]
\centering
\caption{Discrete Schr\"odinger eigenvalues vs.\ Bohr formula $E_n = -\aPDL^2/(2n^2)$ (units $\hb = m_e = 1$, $N = 4000$, $L = 25\,a_0^{\PDL}$).}
\bigskip
\begin{tabular}{cllr}
\toprule
$n$ & $E_n^{\mathrm{disc}}$ & $E_n^{\mathrm{Bohr}} = -\aPDL^2/(2n^2)$ & Deviation \\
\midrule
1 & $-2.6645\times10^{-5}$ & $-2.6631\times10^{-5}$ & $4.8\times10^{-4}$ \\
2 & $-6.6611\times10^{-6}$ & $-6.6578\times10^{-6}$ & $4.9\times10^{-4}$ \\
\bottomrule
\end{tabular}
\label{tab:spectrum}
\end{table}

The residual deviation of $\sim 5\times10^{-4}$ is of purely discretisation origin ($O((\Delta r)^2)$), decreasing with $N$, and is unrelated to the $10^{-4}$ deviation of $\aPDL^{-1} = 137.022$ from CODATA 2022.

\subsection{\PDL\ predictions for the hydrogen atom}

\begin{table}[ht]
\centering
\caption{\PDL\ predictions for hydrogen vs.\ CODATA 2022.}
\bigskip
\begin{tabular}{lllr}
\toprule
Quantity & \PDL\ value & CODATA 2022 & Deviation \\
\midrule
$\aPDL^{-1} = \vp^2\rval^2/(9\mu)$ & $137.022$ & $137.036$ & $1.0\times10^{-4}$ \\
$a_0^{\PDL} = \hb/(m_e c\,\aPDL)$ & $52.934\,\mathrm{pm}$ & $52.918\,\mathrm{pm}$ & $3.0\times10^{-4}$ \\
$E_1^{\PDL} = -m_e c^2\aPDL^2/2$ & $-13.593\,\mathrm{eV}$ & $-13.606\,\mathrm{eV}$ & $1.0\times10^{-4}$ \\
$\mu^* = 1836.152670$ (D28--D30) & & $1836.152673$ & $0.002\,\mathrm{ppm}$ \\
\bottomrule
\end{tabular}
\label{tab:predictions}
\end{table}

All deviations from CODATA trace directly to the same $10^{-4}$ deviation of $\aPDL$. The central advance is that all parameters are now \emph{derived}, not inserted.

% ================================================================
\section{Open problems}
\label{sec:open}

\textbf{OP1 (Algebraic derivation of $b$ without the continuum action argument).} Proposition~\ref{prop:P2} uses the kinetic action $S_{\mathrm{kin}} = m_e(\Delta x)^2/(2\Delta t)$ as an intermediate step. A fully algebraic proof showing directly from the $(A)\wedge(B)$ discrete graph structure that the transition amplitude is purely imaginary with magnitude $\hb/(2m_e)$ in discrete units would constitute the deepest level of the derivation.

\textbf{OP2 (Linearity and unitarity from axioms).} The linear update rule~\eqref{eq:update} is introduced as an effective coarse-grained description. A derivation showing that the \PDL\ logical optimisation functional forces linearity and norm-preservation on $\{\psi_n\}$ would close the last structural gap.

\textbf{OP3 (Multi-body extension).} Theorem~\ref{thm:main} covers one $K_4$ block in the field of one proton. Extension to two electrons (helium) requires specifying how two independent $(A)\wedge(B)$ criteria interact, and whether an effective exchange antisymmetry emerges from the mixed-triangle structure.

\textbf{OP4 (Relativistic extension — Dirac equation).} The $K_4$ pulsation is a two-state object ($s^{(1)}, s^{(2)} = -s^{(1)}$), suggesting a natural spinor representation. The passage from~\eqref{eq:Schrodinger} to the Dirac equation would require restoring the full $SL(2,\mathbb{C})$ structure of the Compton pulsation.

\textbf{OP5 (Born rule from Gleason-type uniqueness).} The complex amplitudes $\psi_n$ are introduced without deriving the Born rule $|\psi_n|^2$. Document~\cite{Gleason_PDL_2026} initiates a Gleason-type approach in the \PDL\ setting; completing it would yield a fully internal derivation of the probability interpretation.

\textbf{OP6 (Connection between $\aPDL$ and $G_{\mathrm{eff}}(N) = \sigma(N)\cdot G_{\PDL}$).} The link between the Schr\"odinger coupling $\aPDL$ and the gravitational coupling of D31 at the level of the 18-step chain complex $\varepsilon_G^{18}$ remains to be established algebraically (D31, OP1).

% ================================================================
\section*{Conclusion}

This paper has derived the non-relativistic Schr\"odinger equation for an electron in the Coulomb field of a \PDL\ proton from the $(A)\wedge(B)$ coupling criterion of D29, the $S_4$ transitivity theorem of D31, and the Compton pulsation of $K_4$ --- without free parameter.

The central advance over the prior sketch documents is the derivation of the transition coefficient $b$. In those documents, $b$ was a phenomenological parameter constrained but not determined by the \PDL\ axioms. Here, condition~(B) --- $P_2 = -P_1$, verified exhaustively over 768 cases with zero counter-example --- forces a $\pi$ phase shift between the two half-cycles of every stable mixed triangle. The coherent sum over one pulsation cycle (equation~\eqref{eq:coherent_sum}) gives a constructive amplitude $\Delta = 2P_1$, and the Compton action $\hb\omega_c = m_e c^2$ then fixes $b$ uniquely as the purely imaginary quantity~\eqref{eq:b}. The $S_4$ transitivity of D31 simultaneously forces the nearest-neighbour symmetric structure of admissible transitions (Proposition~1) and, through the $(A)\wedge(B)$-filtered count $N_{\mathrm{mix}}(r)\propto\Rsurf/r^2$ (Proposition~3), the Coulomb form of the potential with coupling $\aPDL = \mu/\Rsurf^2$ from D29--D30.

The central theorem assembles these four propositions into equation~\eqref{eq:Schrodinger}, resolving open problem~OP6 of D31. Numerical checks confirm: 0/768 counter-examples to the $(A)\wedge(B)$-stability equivalence; orbit matrix of $S_4$ uniformly equal to~4; discrete tridiagonal Hamiltonian converging to the Bohr levels at the $5\times10^{-4}$ level.

% ================================================================
\bibliographystyle{plain}
\bibliography{References_D32}
\end{document}